\chapter{סיכום ומבט לעתיד}

\begin{epigraph}{}
השאלה איננה מה הטכנולוגיה תעשה לנו. השאלה היא מה נבחר לעשות איתה.
\end{epigraph}

\section{מה למדנו}

התחלנו את המסע ביד שנשלחת אל הטלפון עוד לפני שהעיניים נפתחות. בדרך עברנו מהמברק של 1844 אל הסמארטפון של 2007, וראינו איך המידע, שבמשך רוב ההיסטוריה נע במהירות של רגליים, התחיל לנוע במהירות האור.

\begin{sloppypar}
כמה מספרים מסכמים את התמונה. כ-5.87 מיליארד בני אדם השתמשו באינטרנט ב-2024 \parencite{itu2026data}, ובישראל כ-91.1\% מהאוכלוסייה מחוברים \parencite{kemp2025israel}. משתמש טיפוסי מבלה ברשת כ-6 שעות ו-40 דקות ביום \parencite{kemp2024}. בבתים, במפעלים וברחובות פועלים כ-18.5 מיליארד מכשירים מחוברים \parencite{iotanalytics2025}. והמחקר מרמז שבסמארטפון, נקודת המפנה של הרווחה בימי חול מגיעה כבר אחרי כשעתיים ביום \parencite{przybylski2017}.
\end{sloppypar}

אבל הלקח החשוב ביותר איננו מספר. הוא מנגנון. ראינו שהקשב הפך למשאב הנדיר של התקופה \parencite{simon1971}, ושסביבו נבנתה כלכלה שלמה, שהופכת כל גלילה לנתון וכל נתון לתחזית \parencite{zuboff2019}. ראינו שהקשר בין מסכים לרווחה אמיתי, אבל קטן ומורכב, ושהמדענים עצמם עדיין מתווכחים עליו \parencite{orben2019,haidt2024}. וראינו שהאיזון אפשרי, לא דרך התנזרות, אלא דרך בחירה מודעת.

\section{מבט קדימה: המכונה שמדברת}

פגשנו את \en{ChatGPT} כבר בפרק השני: מוצר שהושק בסוף 2022 בלי תרועות, כניסוי \parencite{openai2022}. בדיעבד, זה עשוי להתברר כרגע היסטורי בסדר גודל של המברק.

עד עכשיו, כל טכנולוגיות התקשורת שתיארנו בספר הזה העבירו מילים שבני אדם כתבו. המברק, הטלפון, הרדיו, האינטרנט והרשתות החברתיות היו כולם צינורות. הבינה המלאכותית היא משהו אחר. היא לא רק מעבירה את הסיפורים שלנו. היא מספרת סיפורים משלה.

מה זה יעשה לכלכלת הקשב? אם עד היום אלגוריתמים \textbf{בחרו} איזה תוכן להראות לנו, בקרוב הם עשויים \textbf{ליצור} אותו, במיוחד עבורנו, בכל רגע. מצד אחד, זו הבטחה מרגשת: מורה פרטי לכל ילד, עוזר אישי לכל אדם, ידע שזמין בכל שפה. מצד שני, מערכת שיודעת לשוחח איתנו בחום ובסבלנות אינסופית עלולה להיות מושכת עוד יותר מכל פיד.

כל הכלים שפגשנו בספר הזה רלוונטיים גם כאן. מודל גולדילוקס מזכיר לנו לשאול לא רק "האם", אלא "כמה ומאיזה סוג". המעגל של כלכלת הקשב מזכיר לנו לשאול מי משלם ומי מרוויח. וארגז הכלים מזכיר לנו שגם מול מערכות חכמות מאוד, יש לנו עדיין שליטה על ברירות המחדל של החיים שלנו.

\begin{alertbox}[השאלות החדשות]
\begin{itemize}
  \item כשמכונה כותבת, איך נדע מה אמיתי?
  \item כשעוזר דיגיטלי מכיר אותנו טוב יותר מחברים, למי נפנה ברגע של משבר?
  \item ומי יקבע מה הבינה המלאכותית תגיד לנו, ולמה?
\end{itemize}
\end{alertbox}

\section{שאלה אחרונה}

בני האדם תמיד היו יצורים מחוברים. הציידים-לקטים היו תלויים בשבט שלהם כדי לשרוד. הקשרים האלה הם שהפכו קוף חסר חשיבות מאפריקה לשליט כדור הארץ. במובן הזה, הסמארטפון לא המציא שום דבר חדש. הוא רק לקח את הצורך העתיק ביותר שלנו, הצורך להיות חלק ממשהו, והגביר אותו לעוצמה שלא הכרנו מעולם.

השאלה איננה אם להיות מחוברים. זה כבר הוחלט, ובצדק. השאלה היא למה אנחנו מתחברים, ומה אנחנו מוכנים לתת בתמורה.

\begin{thinkbox}[השאלה שנשארת איתכם]
אם הקשב שלכם הוא המשאב היקר ביותר שיש לכם, ואם יש בעולם מערכות חכמות מאוד שמתחרות עליו בכל רגע, את מי אתם בוחרים לשים במרכז שלו?
\end{thinkbox}

אולי כדאי לחשוב על זה רגע. ואולי, רק הפעם, לעשות את זה בלי הטלפון ביד.
